\section{Base de Datos y Persistencia}

\subsection{Esquema de Base de Datos MySQL}

\subsubsection{Diagrama Entidad-Relación}

\begin{figure}[H]
\centering
\begin{tikzpicture}[
    entity/.style={rectangle, draw, fill=yellow!20, minimum width=3cm, minimum height=0.7cm},
    node distance=2cm
]
    \node[entity] (user) {users};
    \node[entity, below left=of user] (config) {automation\_configs};
    \node[entity, below=of user] (notif) {notifications};
    \node[entity, below right=of user] (report) {reports};
    \node[entity, right=3cm of user] (voice) {voice\_commands};
    
    \draw[-{Latex}] (user) -- node[left,sloped,above] {1:1} (config);
    \draw[-{Latex}] (user) -- node[left] {1:N} (notif);
    \draw[-{Latex}] (user) -- node[right,sloped,above] {1:N} (report);
\end{tikzpicture}
\caption{Diagrama ER del Sistema}
\end{figure}

El sistema utiliza 5 tablas principales con las siguientes relaciones:

\textbf{users} (Tabla central):
\begin{itemize}
    \item 1:1 con automation\_configs
    \item 1:N con notifications
    \item 1:N con reports
\end{itemize}

\textbf{voice\_commands}: Tabla independiente sin relaciones FK

\subsubsection{Tabla: users}

\begin{table}[H]
\centering
\small
\begin{tabular}{|l|l|l|p{4cm}|}
\hline
\textbf{Campo} & \textbf{Tipo} & \textbf{Null} & \textbf{Descripción} \\
\hline
id & bigint & NO & PK, auto\_increment \\
name & varchar(255) & NO & Nombre completo \\
email & varchar(255) & NO & UNIQUE, email para login \\
password & varchar(255) & NO & Hash bcrypt \\
role & varchar(50) & NO & user/admin \\
header\_position & varchar(20) & YES & top/left/right/bottom \\
remember\_token & varchar(100) & YES & Token ``remember me'' \\
email\_verified\_at & timestamp & YES & Verificación email \\
created\_at & timestamp & NO & Fecha creación \\
updated\_at & timestamp & NO & Última actualización \\
\hline
\end{tabular}
\caption{Estructura de tabla users}
\end{table}

\subsubsection{Tabla: automation\_configs}

\begin{table}[H]
\centering
\small
\begin{tabular}{|l|l|p{6cm}|}
\hline
\textbf{Campo} & \textbf{Tipo} & \textbf{Descripción} \\
\hline
id & bigint & PK \\
user\_id & bigint & FK → users.id \\
traffic\_threshold & int & Límite de vehículos para alerta \\
traffic\_minutes & int & Ventana de tiempo (minutos) \\
notify\_high\_traffic & boolean & Activar notificación alto tráfico \\
notify\_camera\_offline & boolean & Activar alerta cámara offline \\
camera\_offline\_minutes & int & Tiempo sin detección para offline \\
report\_frequency & varchar(20) & daily/weekly/monthly/disabled \\
report\_generation\_time & time & Hora de generación automática \\
auto\_generate\_reports & boolean & Activar generación automática \\
xml\_cleanup\_frequency & varchar(20) & Frecuencia de limpieza XML \\
xml\_retention\_days & int & Días de retención de datos \\
\hline
\end{tabular}
\caption{Estructura de automation\_configs}
\end{table}

\subsubsection{Tabla: notifications}

\begin{table}[H]
\centering
\small
\begin{tabular}{|l|l|p{5cm}|}
\hline
\textbf{Campo} & \textbf{Tipo} & \textbf{Descripción} \\
\hline
id & bigint & PK \\
user\_id & bigint & FK → users.id \\
title & varchar(255) & Título de notificación \\
message & text & Mensaje completo \\
type & varchar(50) & info/alert/warning/success \\
priority & varchar(20) & low/medium/high \\
is\_read & boolean & Estado de lectura \\
read\_at & timestamp & Fecha/hora de lectura \\
metadata & json & Datos adicionales \\
created\_at & timestamp & Fecha de creación \\
\hline
\end{tabular}
\caption{Estructura de notifications}
\end{table}

\subsubsection{Tabla: reports}

\begin{table}[H]
\centering
\small
\begin{tabular}{|l|l|p{5cm}|}
\hline
\textbf{Campo} & \textbf{Tipo} & \textbf{Descripción} \\
\hline
id & bigint & PK \\
user\_id & bigint & FK → users.id \\
title & varchar(255) & Título del reporte \\
description & text & Descripción \\
type & varchar(50) & daily/weekly/monthly/custom \\
status & varchar(50) & generated/viewed \\
report\_date & date & Fecha del reporte \\
data & json & Estadísticas generadas \\
filters & json & Filtros aplicados \\
generated\_at &  timestamp & Fecha generación \\
viewed\_at & timestamp & Fecha visualización \\
\hline
\end{tabular}
\caption{Estructura de reports}
\end{table}

\subsubsection{Tabla: voice\_commands}

\begin{table}[H]
\centering
\small
\begin{tabular}{|l|l|p{5cm}|}
\hline
\textbf{Campo} & \textbf{Tipo} & \textbf{Descripción} \\
\hline
id & bigint & PK \\
name & varchar(255) & Nombre del comando \\
trigger & text & Palabras clave (CSV) \\
action & varchar(50) & navigate/export/toggle/custom \\
target & varchar(255) & Ruta o elemento \\
function\_name & varchar(255) & Función JS a ejecutar \\
modules & varchar(255) & Módulos activos (all, modulo1...) \\
enabled & boolean & Estado \\
created\_at & timestamp & Fecha creación \\
\hline
\end{tabular}
\caption{Estructura de voice\_commands}
\end{table}

\newpage

\section{Rutas y API}

\subsection{Rutas Web (routes/web.php)}

\subsubsection{Autenticación}

\begin{table}[H]
\centering
\begin{tabular}{|l|l|l|}
\hline
\textbf{Ruta} & \textbf{Método} & \textbf{Acción} \\
\hline
/login & GET & AuthController@showLogin \\
/login & POST & AuthController@login \\
/logout & POST & AuthController@logout \\
\hline
\end{tabular}
\caption{Rutas de Autenticación}
\end{table}

\subsubsection{Módulos Principales}

\begin{table}[H]
\centering
\begin{tabular}{|l|l|l|}
\hline
\textbf{Ruta} & \textbf{Middleware} & \textbf{Vista/Controller} \\
\hline
/ & - & view('modulo1') \\
/modulo1 & - & view('modulo1') \\
/modulo2 & - & ReporteController@index \\
/modulo3 & isAdmin & ModuloTresController@index \\
/modulo4 & - & view('modulo4') \\
/modulo5 & isAdmin & UserController@index \\
\hline
\end{tabular}
\caption{Rutas de Módulos}
\end{table}

\subsubsection{APIs}

\begin{table}[H]
\centering
\small
\begin{tabular}{|l|l|p{5cm}|}
\hline
\textbf{Endpoint} & \textbf{Método} & \textbf{Descripción} \\
\hline
/api/vehicle-monitor/\{cameraId\} & GET & Estadísticas del detector \\
\hline
/api/voice-commands & GET & Listar comandos de voz \\
/api/voice-commands & POST & Crear comando \\
/api/voice-commands/\{id\} & PUT & Actualizar comando \\
/api/voice-commands/\{id\}/toggle & POST & Activar/desactivar \\
\hline
/api/notifications & GET & Listar notificaciones (admin) \\
/api/notifications/unread & GET & Contador no leídas \\
/api/notifications/\{id\}/read & POST & Marcar como leída \\
\hline
/api/reports & GET & Listar reportes (admin) \\
/api/reports/generate & POST & Generar reporte \\
/api/reports/\{id\} & GET & Ver reporte específico \\
\hline
\end{tabular}
\caption{Endpoints API}
\end{table}

\newpage

\section{Patrones de Diseño Implementados}

\subsection{MVC (Model-View-Controller)}

Patrón fundamental de Laravel, separa lógica de negocio, presentación y datos.

\begin{figure}[H]
\centering
\begin{tikzpicture}[
    component/.style={rectangle, draw, fill=blue!20, minimum width=2.5cm, minimum height=1cm},
    node distance=3cm
]
    \node[component] (view) {VIEW\\Blade};
    \node[component, right=of view] (controller) {CONTROLLER\\PHP};
    \node[component, below=2cm of controller, xshift=-1.5cm] (model) {MODEL\\Eloquent};
    
    \draw[-{Latex}] (view) -- node[above] {Request} (controller);
    \draw[-{Latex}] (controller) -- node[below,sloped] {Response} (view);
    \draw[-{Latex}] (controller) -- node[right] {Query} (model);
    \draw[-{Latex}] (model) -- node[left] {Data} (controller);
\end{tikzpicture}
\caption{Arquitectura MVC del Sistema}
\end{figure}

\textbf{Controllers}: AuthController, VehicleMonitorController, VoiceCommandController, etc.

\textbf{Models}: User, Notification, Report, AutomationConfig

\textbf{Views}: modulo1.blade.php, modulo2.blade.php, modulo3.blade.php, modulo4.blade.php, modulo5.blade.php

\subsection{Repository/Service Pattern}

Abstrae la lógica de acceso a datos y lógica de negocio compleja.

\textbf{Implementación}:

\begin{itemize}
    \item \texttt{NotificationService}: Centraliza lógica de notificaciones
    \item \texttt{ReportService}: Generación compleja de reportes desde XML
    \item \texttt{NotificationFacade}: Facade pattern para simplificar acceso
    \item \texttt{ReportFacade}: Facade pattern para reportes
\end{itemize}

\subsection{Observer Pattern}

Eloquent de Laravel implementa observers para eventos del modelo:
\begin{itemize}
    \item creating, created
    \item updating, updated
    \item deleting, deleted
\end{itemize}

\subsection{Factory Pattern}

Usado en seeders y testing para generar datos de prueba mediante UserFactory.

\subsection{Strategy Pattern}

Implementado en el sistema de detección multi-cámara:

\begin{figure}[H]
\centering
\begin{tikzpicture}[
    box/.style={rectangle, draw, fill=orange!20, minimum width=3cm, minimum height=0.8cm},
    node distance=1.5cm
]
    \node[box] (context) {camera\_states\\Contexto};
    \node[box, below=of context] (interface) {CameraStrategy\\Interfaz};
    \node[box, below left=of interface] (oracle) {worker\_oracle()\\WebSocket};
    \node[box, below right=of interface] (skyline) {worker\_skyline()\\HLS};
    
    \draw[-{Latex}] (context) -- (interface);
    \draw[-{Latex}, dashed] (interface) -- (oracle);
    \draw[-{Latex}, dashed] (interface) -- (skyline);
\end{tikzpicture}
\caption{Strategy Pattern en Detección Multi-Cámara}
\end{figure}

\textbf{Contexto}: camera\_states (gestión de estado por cámara)

\textbf{Estrategias}:
\begin{itemize}
    \item worker\_oracle(): WebSocket Strategy para CAM\_001
    \item worker\_skyline(): HLS Strategy para CAM\_002
\end{itemize}

Cada estrategia implementa su propia lógica de conexión y procesamiento de frames.

\newpage

\section{Flujos de Trabajo Completos}

\subsection{Flujo: Detección de Vehículo}

\begin{enumerate}
    \item Cámara captura frame
    \item Worker recibe frame
    \item YOLO detect\_vehicles() - detección
    \item match\_vehicle\_to\_tracked() - tracking
    \item get\_dominant\_color() - análisis de color
    \item ¿Es nuevo vehículo? Si → save\_detection\_to\_xml()
    \item WebSocket broadcast a clientes conectados
    \item Canvas en navegador dibuja frame
\end{enumerate}

\subsection{Flujo: Generación de Reporte}

\begin{enumerate}
    \item Usuario accede a Módulo 2 o API solicita generación
    \item ReportController recibe request con filtros (tipo, fechas)
    \item ReportService::generate() es invocado
    \item generateReportData() lee vehiculos\_db.xml
    \item Itera sobre \textless{}deteccion\textgreater{} aplicando filtros
    \item Acumula estadísticas (by\_type, by\_hour, by\_day, by\_color)
    \item Calcula promedio de confianza, hora pico, día más ocupado
    \item Crea registro en tabla reports con data JSON
    \item Retorna Report model con estadísticas completas
    \item Usuario visualiza o descarga reporte
\end{enumerate}

\subsection{Flujo: Comando de Voz}

\begin{enumerate}
    \item Usuario habla al micrófono
    \item Vosk procesa audio
    \item SocketIO emite 'voice\_command'
    \item JavaScript recibe evento
    \item findMatchingCommand(text) busca coincidencias
    \item executeCommand() ejecuta acción
    \item Acción: navigate, export, o toggle según configuración
\end{enumerate}

\newpage

\section{Seguridad}

\subsection{Autenticación}

\begin{itemize}
    \item \textbf{Hashing de Contraseñas}: Bcrypt con factor de costo configurable
    \item \textbf{Sesiones}: Laravel session management con regeneración en login
    \item \textbf{Remember Token}: Token de 100 caracteres para ``recordarme''
    \item \textbf{CSRF Protection}: Tokens CSRF en todos los formularios
\end{itemize}

\subsection{Autorización}

\begin{itemize}
    \item \textbf{Middleware isAdmin}: Protege rutas administrativas (modulo3, modulo5)
    \item \textbf{Roles}: user / admin
    \item \textbf{Policy-based}: Verificación de propiedad de recursos
\end{itemize}

\subsection{Validación de Datos}

\begin{itemize}
    \item \textbf{Form Requests}: Validación automática con reglas Laravel
    \item \textbf{Sanitización}: Escape de HTML en outputs Blade
    \item \textbf{SQL Injection}: Protección mediante Eloquent ORM
\end{itemize}

\newpage

\section{Configuración y Deployment}

\subsection{Requisitos del Sistema}

\begin{table}[H]
\centering
\begin{tabular}{|l|l|}
\hline
\textbf{Componente} & \textbf{Versión Mínima} \\
\hline
PHP & 8.2+ \\
Laravel & 11.x \\
Python & 3.9+ \\
MySQL & 8.0+ \\
Node.js & 18+ (para npm) \\
Composer & 2.x \\
\hline
\end{tabular}
\caption{Requisitos del Sistema}
\end{table}

\subsection{Puertos y Servicios}

\begin{table}[H]
\centering
\begin{tabular}{|l|l|l|}
\hline
\textbf{Servicio} & \textbf{Puerto} & \textbf{Comando} \\
\hline
Laravel (Apache) & 80 & php artisan serve \\
MySQL & 3306 & mysqld \\
Detector Python & 8080 & python vehiculo\_detector.py \\
Voice Server & 5001 & python voice\_server.py \\
\hline
\end{tabular}
\caption{Puertos y Servicios}
\end{table}

\newpage

\section{Conclusiones}

\subsection{Logros del Sistema}

\begin{enumerate}
    \item \textbf{Detección en Tiempo Real}: YOLOv8 con precisión superior al 85\%
    \item \textbf{Multi-Cámara}: Soporte para 2+ cámaras simultáneas con tracking independiente
    \item \textbf{WebSocket Streaming}: Latencia inferior a 100ms
    \item \textbf{Comandos de Voz}: Reconocimiento en español con Vosk
    \item \textbf{Automatización}: Sistema completo de notificaciones y reportes
    \item \textbf{Escalabilidad}: Arquitectura modular preparada para crecer
\end{enumerate}

\subsection{Características Técnicas Destacadas}

\begin{itemize}
    \item \textbf{Persistencia Dual}: MySQL para datos estructurados, XML para detecciones
    \item \textbf{Tracking Único}: Algoritmo basado en distancia euclidiana
    \item \textbf{Detección de Color}: Análisis HSV para clasificación cromática
    \item \textbf{API RESTful}: Endpoints bien documentados y consistentes
    \item \textbf{Responsive UI}: TailwindCSS para diseño adaptable
\end{itemize}

\end{document}
